\section{A. Groundwater accounts}\label{sec:a}

\begin{inshort}
For each regency and city of Central Java, module A puts three numbers side by side: what was declared,
what our validated method estimates, and what the difference is worth under the official tariff. The
tariff side is fully official. The gap is not known yet, and the app says so with a range.
\end{inshort}

\subsection{Declared use}

Dinas ESDM Jateng publishes the declared abstraction volume and its NPA for every regency and city, as XLSX
files for 2020, 2021, 2022 and 2023 \cite{esdmdata}. Kota Semarang declared 10,069,814\,\mthree\ in 2023 \Off.
A second data set gives use as a share of availability, 2019--2023 \cite{esdmpct}. Module A can therefore
cover all 35 regencies and cities, not only Semarang.

The city's own records show a longer series of registered abstraction \cite{perdasmg}:

\begin{center}
\small
\begin{tabular}{@{}lrrrrrrrr@{}}
\toprule
Year & 1990 & 1995 & 2000 & 2016 & 2017 & 2018 & 2019 & 2023\\
\midrule
Wells & -- & 300 & 320 & 1,050 & -- & -- & -- & --\\
Volume (\Mm) & 23 & 27 & 38 & 13.1 & 15.5 & 27.7 & 13.8 & 10.1\\
\bottomrule
\end{tabular}
\end{center}

\begin{caution}
The series jumps from 13.8 to 27.7\Mm\ between neighbouring years. The declared numbers are themselves
noisy, so a fall from 38 (2000) to 10 (2023) mixes three things: real decline, under-reporting and wells
that were never registered.
\end{caution}

\subsection{Who pumps}

The same document says 96\,\% of registered abstraction is industrial. PDAM (the city water company)
supplied 62\,\% of household demand and distributed 51.66\Mm\ in 2022: 41.63 to households, 4.35 to
commerce and \textbf{only 0.22 to industry} \Off\ \cite{perdasmg}. Industry therefore runs almost entirely
on its own groundwater. The text also says about 19\,\% of supply comes from deep groundwater; the wording
does not make clear whether this is PDAM's own wells.

\subsection{The official price}\label{sec:npa}

Pergub Jateng 7/2025 (in force 18 February 2025) sets the groundwater price, NPA, as \cite{pergub7}:
\[
\text{NPA} = \text{HAB} \times \text{BAT},\qquad \text{BAT} = 0.6\,S + 0.4\,P .
\]
HAB is the base price of water: Rp 2,561 (Wilayah A), 2,795 (B) and 2,824 (C) per \mthree. Wilayah C is
Kota Semarang, Kab.\ Demak, Kab.\ Kudus, Kab.\ Pekalongan and Kota Surakarta. $S$ is the resource
component (16, 9, 4 or 1, by water quality and whether an alternative source exists). $P$ is the use group:
Kelompok 1 sells water as a product, 2--4 are high-, medium- and low-risk businesses, 5 is social,
government and public water supply. Five volume tiers run from 0--50 to more than 2,500\,\mthree. The pages
read do not state the period the tiers refer to.

\begin{table}[H]
\centering\small
\caption{NPA in Wilayah C, good quality with an alternative source (annex table 3a), and the tax at 20\,\%.
All values \Off\ except the tax column \Calc.}\label{tab:npa}
\begin{tabular}{@{}lrrr@{}}
\toprule
\textbf{Use group} & \textbf{NPA, 0--50\,\mthree} & \textbf{NPA, $>$2,500\,\mthree} & \textbf{Tax, top tier}\\
 & Rp/\mthree & Rp/\mthree & Rp/\mthree\\
\midrule
K1 water as a product & 37,277 & 78,575 & 15,715\\
K2 high-risk business & 35,018 & 67,143 & 13,429\\
K3 medium-risk business & 32,758 & 55,701 & 11,140\\
K4 low-risk business & 30,499 & 44,269 & 8,854\\
K5 social, government, public supply & 28,240 & 32,826 & 6,565\\
\bottomrule
\end{tabular}
\end{table}

The annex prints the 40\,\% $P$ column for Wilayah C as 0.40--2.02 for every group, which looks like a typo;
the BAT totals match Wilayah A, so the NPA column is used as printed.

\textbf{Tax rate.} Kota Semarang taxes groundwater (PAT) at 20\,\% of NPA under Perda 10/2023, amended by
Perda 4/2025 \cite{ddtcsmg,perda10}. 20\,\% is also the maximum under UU 1/2022 \cite{terc}. This comes
from a tax news service and a policy brief; the Perda clause itself was not read \Wk.

\textbf{Who is exempt.} Basic household use is reported as exempt, by one secondary source \cite{enforcea}
\Wk. The UU 1/2022 article must be read before households are left out of the rupiah figure.

\subsection{Method \M{1}: declared vs estimated}

The estimate comes from two earlier results:
\begin{itemize}
\item the demand model of file 07 at 4\,km cells, which gives Kota Semarang 25--48\Mmyr\ (median 36.5) but
  was never validated \cite{exp07};
\item the validated split of file 09 at 16\,km, which gives the west and central cell 27.9 / 43.8 /
  54.7\Mmyr\ (P10 / P50 / P90) \cite{exp09} \Tst.
\end{itemize}

The taxable gap is then
\[
\text{gap} = \text{estimated total} - \text{household wells (exempt)} - \text{public supply wells} - \text{declared}.
\]

\begin{table}[htbp]
\centering\small
\caption{The gap for Kota Semarang with today's inputs.}\label{tab:gap}
\begin{tabularx}{\linewidth}{@{}L{48mm}L{26mm}Y l@{}}
\toprule
\textbf{Term} & \textbf{Value, \Mmyr} & \textbf{Basis} & \textbf{Label}\\
\midrule
Estimated total & 25--55 & Files 07 and 09 & \Tst\ / \Calc\\
Household wells outside PDAM & 0--25 & PDAM households $41.63\times38/62$, an upper bound that includes shallow wells and springs & \Calc\\
Public supply wells & 0--10 & About 19\,\% of PDAM supply, if these wells are not already in the declared volume & \Wk\\
Declared (2023) & 10.1 & ESDM open data & \Off\\
\midrule
\textbf{Taxable gap} & \textbf{below 0 to about 45} & Low: $25-25-10-10$; high: $55-0-0-10$ & \Calc\\
\bottomrule
\end{tabularx}
\end{table}

\begin{caution}
With today's inputs we cannot say that there \emph{is} a gap. The household term is the largest unknown.
The app shows this range and the terms that make it up; it does not print a single ``lost revenue''
number. The range narrows when \V{1} is run and when household well use is measured or requested.
\end{caution}

\subsection{Method \M{2}: pricing}

Pricing is plain arithmetic on official numbers: tax per \mthree\ $=0.20\times$ NPA \Calc. For business use in
the top tier (K4 to K2), each 1\Mm\ of undeclared pumping is \textbf{Rp 8.9--13.4 billion per year} in tax.
Pricing the declared 10.07\Mm\ at the K3--K2 top tier gives Rp 112--135 billion per year; this is the
\emph{size} of the tax base, not a measured revenue.

The app shows a slider for ``undeclared business use'' \Scn, starting at zero, and prices whatever the user
sets. Earlier drafts printed a scenario of 5--15\Mm\ (Rp 45--200 billion); this document drops it as a
finding because nothing measured backs it.

\subsection{The first check, \V{1}: billing}\label{sec:v1}

Get the realised groundwater tax for Kota Semarang, 2020--2025, from the budget realisation report (APBD)
or the mayor's accountability report (LKPJ). Divide it by the declared volume of the same year. The
implied price per \mthree\ should fall inside $0.20\times$ the NPA range of groups 2--4.
\begin{itemize}
\item Inside the range: billing matches the declarations; any gap is about undeclared volume.
\item Below the range: tax is leaking at the billing stage, before any unregistered well is counted. This is
  the strongest single number for the pitch, and it needs no satellite data.
\end{itemize}

\begin{caution}
Pergub 7/2025 applies from February 2025. For 2020--2024 the previous NPA regulation applies, and we have
not obtained it yet (\R{1}).
\end{caution}
